% !TEX root = ../ComputationalOTFnT.tex

\chapter{動的定式化}
\label{c-dynamic}

この章では、コストが測地距離の二乗である場合の最適輸送の測地線的（動的とも呼ばれる）視点を提示する。これは、二つの測度間の最適輸送を、測度の空間における全長を最小化する曲線として記述するものである。動的な視点は最適輸送の直感的な代替解釈を提供し、流体力学との関連を引き出すだけでなく、二つの密度間の補間を行う際に低次元での OT 計算のための効率的な数値ツールをもたらす。この手法の欠点は、大規模な疎な測度には適用できず、測地距離の二乗コストを用いた正則な領域上の低次元（空間を格子化する必要があるため）でのみ機能することである。

この章では、時刻 $t=0$ の一つの測度から出発して時刻 $t=1$ で別の測度に到達するという考え方に基づき、$(\al,\be)$ の代わりに $(\al_0,\al_1)$ という表記を用いる。

%%%%%%%%%%%%%%%%%%%%%%%%%%%%%%%%%%%%%%%%%
\section{連続定式化}

$\X=\Y=\RR^d$ かつ $c(x,y)=\norm{x-y}^2$ の場合、\eqref{eq-mk-generic}で定義される最適輸送距離 $\Wass_2^2(\al,\be) = \MK_\c(\al,\be)$ は、二つの測度間の最短経路 $(\al_t)_{t=0}^1$ を探索することによって計算できる。
%
この経路は、各時刻で定義されるベクトル場 $\speed_t$ を用いて測度を移流させることで記述される。ベクトル場 $\speed_t$ と経路 $\alpha_t$ は質量保存の公式を満たさなければならず、その結果として
\eql{\label{eq-cons-bb-ncvx}
	\pd{\al_t}{t} + \diverg(\al_t \speed_t) = 0
	\qandq
	\al_{t=0}=\al_0,
	\al_{t=1}=\al_1,
}
が得られる。ここで上式は $\RR^d$ 上の超関数の意味で理解されるべきである。このようなベクトル場の微小長さは、測度 $\al_t$ に付随する $L^2$ ノルムを用いて測定され、次のように定義される。
\eq{\norm{\speed_t}_{L^2(\al_t)} = \left( \int_{\RR^\dim} \norm{\speed_t(x)}^2 \d\al_t(x) \right)^{1/2}.}
この定義は、\citet{benamou2000computational}によって最初に導入された $\Wass_2$ の以下の最短経路再定式化をもたらす：
\eql{\label{eq-bb-non-cvx}
	\Wass_2^2(\al_0,\al_1) = \umin{ (\al_t,\speed_t)_t \text{ sat. \eqref{eq-cons-bb-ncvx}} }
		\int_0^1 \int_{\RR^d}   \norm{\speed_t(x)}^2 \d\al_t(x) \d t,
}
ここで $\al_t$ はスカラー値測度、$\speed_t$ はベクトル値測度である。図~\ref{fig-displacement}にこのような測度の経路の二つの例を示す。


\newcommand{\MyFigDispl}[1]{\includegraphics[width=.19\linewidth]{displacement/#1}}
\begin{figure}[h!]
\centering
\begin{tabular}{@{}c@{\hspace{1mm}}c@{\hspace{1mm}}c@{\hspace{1mm}}c@{\hspace{1mm}}c@{}}
\MyFigDispl{interp-dens-1} &
\MyFigDispl{interp-dens-2} &
\MyFigDispl{interp-dens-3} &
\MyFigDispl{interp-dens-4} &
\MyFigDispl{interp-dens-5} \\
\MyFigDispl{interp-points-1} &
\MyFigDispl{interp-points-2} &
\MyFigDispl{interp-points-3} &
\MyFigDispl{interp-points-4} &
\MyFigDispl{interp-points-5} \\
$\al_0$ & $\al_{1/4}$ & $\al_{1/2}$ & $\al_{3/4}$ & $\al_{1}$
\end{tabular}
\caption{\label{fig-displacement}
\eqref{eq-bb-non-cvx}を満たす変位補間 $\al_t$。
上段：Lebesgue測度に関する密度を持つ二つの測度 $(\al_0,\al_1)$ に対する場合。
下段：同数の点を持つ二つの離散経験測度に対する場合。
}
\end{figure}


定式化~\eqref{eq-bb-non-cvx}は、制約~\eqref{eq-cons-bb-ncvx}が積 $\al_t \speed_t$ を含むため、変数 $(\al_t,\speed_t)_t$ に関して非凸な定式化である。ベクトル値測度（しばしば「運動量」と呼ばれる）
\eq{
	\Moment_t \eqdef \al_t \speed_t,
}
を導入すると、\citeauthor{benamou2000computational}は画期的な論文~\citeyearpar{benamou2000computational}において、次のように書くことで変数 $(\al_t,\Moment_t)_t$ に関して凸になることを示した：
\eql{\label{eq-bb-convex}
	\Wass_2^2(\al_0,\al_1) = \umin{ (\al_t,\Moment_t)_t  \in \BBConstr(\al_0,\al_1)  }
		\int_0^1 \int_{\RR^d} \theta(\al_t(x),\Moment_t(x)) \d x \d t,
}
ここで制約集合を次のように定義する：
\eql{\label{eq-mass-conservation}
	\BBConstr(\al_0,\al_1) \eqdef
	\enscond{ (\al_t,\Moment_t) }{
		\pd{\al_t}{t} + \diverg(\Moment_t) = 0,
		\al_{t=0}=\al_0,
		\al_{t=1}=\al_1
	},
}
また $\theta :  \rightarrow \RR^+ \cup \{+\infty\}$ は以下の下半連続凸関数である：
\eql{\label{eq-theta-dynamic}
	\foralls (a,b) \in \RR_+ \times \RR^\dim, \quad
	\th(a,b) = \choice{
		\frac{\norm{b}^2}{a} \qifq a>0, \\
		0 \qifq (a,b) = 0, \\
		+\infty \quad \text{otherwise}.
	}
}
この定義は一見複雑に見えるかもしれないが、$\al_t(x)=0$ のとき運動量 $\Moment_t(x)$ が消滅すべきことを課すために不可欠である。
%
また、\eqref{eq-bb-convex}は測度 $(\al_t,\Moment_t)$ が密度関数であるかのように非形式的に書かれているが、$\th$ は1次斉次関数であるため（したがって測度がLebesgue測度に関する密度を持たない場合でも定義される）、密度から関数への一意的な拡張が可能であり、これは許容される。

%%%%%%%%%%%%%%%%%%%
%\begin{rem11}{Links with McCann's interpolation}\label{rem-displacement}
\begin{rem11}{McCannの補間との関連}{rem-displacement}
（式~\eqref{eq-brenier-map}の場合のように）$\T_\sharp \al_0=\al_1$ を満たす最適Monge写像 $\T : \RR^\dim \rightarrow \RR^\dim$ が存在する場合、$\al_t$ はMcCannの補間
\eql{\label{eq-mc-cann-interp}
	\al_t = ((1-t) \Id + t \T)_{\sharp} \al_0
}
に等しくなる。
1次元の場合、注意~\ref{rem-1d-ot-generic}を用いると、この補間は次の関係式を用いて計算できる：
\eql{\label{eq-displacement-1d-cumul}
	\cumul{\al_t}^{-1} = (1-t)\cumul{\al_0}^{-1} + t \cumul{\al_1}^{-1};
}
図~\ref{fig-1d-ot}を参照されたい。Wasserstein空間のRiemann幾何学の詳細なレビューについては~\citet{gangbo1996geometry}を参照されたい。
%
最適カップリング $\pi$ がMonge写像の台上にあるとは限らず「単に」最適カップリングが存在する場合には、この補間を次のように計算できる：
\eql{\label{eq-displ-discr}
	\al_t = P_{t\sharp} \pi
	\qwhereq
	P_t : (x,y) \in \RR^\dim \times \RR^\dim \mapsto (1-t) x + t y.
}
例えば、離散的な設定~\eqref{eq-pair-discr}において、$\P$ を~\eqref{eq-mk-discr}の解とすると、補間は次のように定義される：
\eql{\label{eq-mccann-discrete}
	\al_t = \sum_{i,j} \P_{i,j} \de_{ (1-t) x_i + t y_j }.
}
このような補間は通常 $n+m-1$ 個の点上に台を持ち、これは $\P$ の非零要素の最大数である。
%
図~\ref{fig-displacement-emp-weight}に離散測度の変位補間の二つの例を示す。
%
この構成は、$P_t$ を測地線経路に沿った補間に置き換えることで、測地空間 $\Xx$ に一般化できる。
%
McCannの補間は、例えばコンピュータグラフィクスにおける色、形状、照明の補間など、多くの応用がある~\citep{Bonneel-displacement}。
%\end{rem1}
\end{rem11}
%%%%%%%%%%%%%%%%%%%



\newcommand{\MyFigDisplEmp}[1]{\includegraphics[width=.16\linewidth]{interpolation-discrete/empirical/interp-#1}}
\newcommand{\MyFigDisplWei}[1]{\includegraphics[width=.16\linewidth]{interpolation-discrete/weighted/interp-#1}}
\begin{figure}[h!]
\centering
\begin{tabular}{@{}c@{}c@{}c@{}c@{}c@{}c@{}}
\MyFigDisplEmp{1}&
\MyFigDisplEmp{2}&
\MyFigDisplEmp{3}&
\MyFigDisplEmp{4}&
\MyFigDisplEmp{5}&
\MyFigDisplEmp{6}\\
\MyFigDisplWei{1}&
\MyFigDisplWei{2}&
\MyFigDisplWei{3}&
\MyFigDisplWei{4}&
\MyFigDisplWei{5}&
\MyFigDisplWei{6}\\
$\al_0$ & $\al_{1/5}$ & $\al_{2/5}$ & $\al_{3/5}$ & $\al_{4/5}$ & $\al_{1}$
\end{tabular}
\caption{\label{fig-displacement-emp-weight}
離散測度の変位補間~\eqref{eq-displ-discr}の比較。
上段：点群（同数の点を持つ経験測度 $(\al_0,\al_1)$）。
下段：同様であるが重みが変化する場合。$0<t<1$ において、上段の例は $N$ 個の点を持つ経験測度の補間 $\mu_t$ に対応し、下段は $2N-1$ 個の点上に台を持つ測度を定義する。
}
\end{figure}

%%%%%%%%%%%%%%%%%%%%%%%%%%%%%%%%%%%%%%%%%
\section{一様スタガード格子上の離散化}

簡単のため、$\dim=2$ 次元の場合に数値スキームを記述する。より高次元への拡張は直接的である。ここでは~\citet{FPapPeyOud13}によって導入された離散化手法に従い、流体力学で一般的に用いられるスタガード格子技法に触発されたものである。
%
時間を $t_k = k/T \in [0,1]$ と離散化し、空間は点 $x_i = (i_1/n_1,i_2/n_2) \in X = [0,1]^2$ で一様に離散化されるものとする。
%
スタガード格子表現を用いるため、$\al_t$ は時間方向の半格子点に対応する $\a \in \RR^{(T+1) \times n_1 \times n_2}$ で表現され、一方 $\Moment$ は $\MomentD=(\MomentD_1,\MomentD_2)$ で表現される。ここで $\MomentD_1 \in \RR^{T \times (n_1+1) \times n_2}$ と $\MomentD_1 \in \RR^{T \times n_1 \times (n_2+1)}$ は各空間方向の半格子点に格納される。
%
この表現を用いると、$(k,i_1,i_2) \in \range{1,T} \times \range{1,n_1} \times \range{1,n_2}$ に対して、時間微分は
\eq{
	(\partial_t \a)_{k,i} \eqdef \a_{k+1,i}-\a_{k,i}
}
と計算され、空間発散は
\eql{\label{eq-space-diverg-disc}
	\diverg(\MomentD)_{k,i} \eqdef
			\MomentD_{k,i_1+1,i_2}^1-\MomentD_{k,i_1,i_2}^1
			+
			\MomentD_{k,i_1,i_2+1}^2-\MomentD_{k,i_1,i_2}^2,
}
%(assuming $0$ whenever these quantities fall out of the grid)
と計算される。これらはいずれも格子点上で定義され、$\RR^{T \times n_1 \times n_2}$ の配列を形成する。

最適化される汎関数を評価するためには、半格子点から格子点への補間演算子が必要であり、すべての $(k,i_1,i_2) \in \range{1,T} \times \range{1,n_1} \times \range{1,n_2}$ に対して、
\eq{
	\InterpBB_a( \a )_{k,i} \eqdef \InterpBB(\a_{k+1,i},\a_{k,i}),
}
\eq{
	\InterpBB_\Moment( \MomentD )_{k,i} \eqdef (
		\InterpBB(\MomentD_{k,i_1+1,i_2}^1,\MomentD_{k,i_1,i_2}^1),
		\InterpBB(\MomentD_{k,i_1,i_2+1}^2,\MomentD_{k,i_1,i_2}^2)
	 ).
}
最も単純な選択は線形演算子 $\InterpBB(r,s)=\frac{r+s}{2}$ を用いることであり、以下ではこれを考える。\eqref{eq-bb-convex}の離散的な対応物は次のようになる：
\eql{\label{eq-bb-discrete}
	\umin{(\a,\MomentD) \in \BBConstrD(\a_0,\a_1)}
		\Theta(\InterpBB_a(\a),\InterpBB_\Moment( \MomentD )),
}
\eq{
		\qwhereq
		\Theta(\tilde\a,\tilde\MomentD)
		\eqdef \sum_{k=1}^{T} \sum_{i_1=1}^{n_1} \sum_{i_2=1}^{n_2}
		\th( \tilde\a_{k,i}, \tilde\MomentD_{k,i} ),
}
ここで制約は
\eq{
	\BBConstrD(\a_0,\a_1) \eqdef
	\enscond{ (\a, \MomentD)
	}{
		\partial_t \a + \diverg(\MomentD) = 0,
		(\a_{0,\cdot},\a_{T,\cdot}) = (\a_0,\a_1)
	},
}
と表される。ここで $\a \in \RR^{(T+1) \times n_1 \times n_2}$、$\MomentD=(\MomentD_1,\MomentD_2)$ で $\MomentD_1 \in \RR^{T \times (n_1+1) \times n_2}$、$\MomentD_2 \in \RR^{T \times n_1 \times (n_2+1)}$ である。
%
図~\ref{fig-dynamic-maze}にこの離散化スキームを用いて近似された発展 $(\al_t)_t$ の例を示す。

%%%%
\begin{rem}[グラフ上の動的定式化]
$\X$ がグラフで $c(x,y)=d_\X(x,y)^2$ が測地距離の二乗である場合、一様格子離散化~\eqref{eq-space-diverg-disc}の代わりにグラフ構造に対応する離散発散を用いた忠実な離散化手法を導出することが可能である。対応する動的Wasserstein距離に対して熱方程式が勾配流構造を持つことを保証するために（勾配流の詳細については~\S\ref{sec-grad-flows}を参照）、\citet{Maas2011}、その後~\citet{MielkeCVPDE}は対数平均 $\InterpBB(r,s)$ の使用を提案した（\citep{solomon2016continuous,ChowHuangLiZhou2012,chow2017entropy,chow2017discrete}も参照）。
\end{rem}
%%%%

%%%%%%%%%%%%%%%%%%%%%%%%%%%%%%%%%%%%%%%%%
\section{近接作用素ソルバー}
\label{sec-prox-solvers}

離散化された動的OT問題~\eqref{eq-bb-discrete}は、アフィン制約の下での非平滑最適化問題の最小化を要求するため、解くことが困難である。実際、関数 $\th$ は凸であるが、質量が消滅する測度 $\a_{k,i}$ に対しては非平滑である。コンパクトな台を持つ二つの入力測度 $(\a_0,\a_1)$ 間の補間を行う場合、補間された測度 $(\a_{k})_{k=1}^T$ の質量も消滅することが一般に期待され、最適化過程の困難な部分はまさにこの台の発展を追跡することにある。特に、標準的な滑らかな最適化技法を用いることはできない。

\eqref{eq-bb-discrete}を二次錐計画として再定式化する方法はいくつかあり、双対問題を考えるか、あるいは単に汎関数 $\th( \a_{k,i}, \MomentD_{k,i} )$ を制約下の線形関数に置き換えることで実現される：
\eq{
	\Theta(\tilde\a,\tilde\MomentD) = \umin{\tilde \zD} \enscond{
		\sum_{k,i} \tilde \zD_{k,i}
	}{
		\foralls (k,i), (\zD_{k,i},\tilde\a_{k,i},\tilde\MomentD_{i,j}) \in \Ll
	},
}
これには追加変数 $\tilde \zD$ の導入が必要となる。
%
ここで $\Ll \eqdef \ensconds{ (z,a,J) \in \RR \times \RR^+ \times \RR^\dim }{ \norm{J}^2 \leq z a }$ は回転ローレンツ二次錐である。
%
この追加変数により、二次錐計画の標準的な内点法ソルバー~\citep{nesterov1994interior}を用いて離散化問題を解くことが可能となる。これらのソルバーは高速な収束率を持ち、高精度の解を計算できる。しかし、各反復は計算コストが高く、離散化点の数に比例する次元の線形系の求解を必要とする。そのため、画像処理応用で遭遇する大規模多次元問題には適用できない。

これらの高精度ソルバーの代替手段として、低精度の一次手法がある。これらは~\eqref{eq-bb-discrete}のような非平滑だが高度に構造化された問題に適している。この種のソルバーは新しいものではないが、数値精度が主要な目標ではないこれらの応用に最適であるため、画像処理や機械学習の分野で最近再び注目を集めている。
%
これらのソルバーとその大規模応用への使用に関する詳細な説明については、例えばモノグラフ~\citep{BauschkeCombettes11}を参照されたい。
%
ここでは特定のソルバーに集中するが、もちろん他の多くのソルバーも使用可能であり、動的OTに対するいくつかのアプローチの研究については~\citep{FPapPeyOud13}を参照されたい。動的OTに一次スキームを用いるというアイデアは、最初に~\citet{benamou2000computational}によって提案されたことに注意されたい。

DRアルゴリズム~\citep{Lions-Mercier-DR}は、次の形式の非平滑構造化問題を解くために特別に設計されている：
\eql{\label{eq-splitting-opt}
	\umin{x \in \Hh} F(x) + G(x),
}
ここで $\Hh$ はEuclid空間であり、$F, G : \Hh \rightarrow \RR \cup \{+\infty\}$ は二つの閉凸関数で、いわゆる近接作用素
\eql{\label{eq-prox-eucl}
	\foralls x \in \Hh, \quad
	\Prox_{\tau F}(x) \eqdef \uargmin{x' \in \Hh} \frac{1}{2}\norm{x-x'}^2 + \tau F(x)
}
がパラメータ $\tau>0$ に対して「容易に」（閉形式で、または高速に収束するスキームを用いて）計算できるものである。これはEuclid計量に対する近接写像に対応し、$\norm{x-x'}^2$ の代わりにより一般的なBregmanダイバージェンスを用いた定義に拡張できることに注意されたい。$\KLD$ ダイバージェンスを用いた例については~\eqref{eq-prox-kl}を参照されたい。
%
DRアルゴリズムの反復は、初期化 $(x^{(0)}, w^{(0)})\in \Hh^2$ を用いて列 $(\it{x}, \it{w}) \in \Hh^2$ を次のように定義する：
\begin{equation}\label{eq-dr-iters}
\begin{aligned}
	\itt{w} &\eqdef  \it{w} + \al (\Prox_{\ga F} (2\it{x}-\it{w})-\it{x}),\\
	\itt{x} &\eqdef  \Prox_{\ga G}(\itt{w}).
\end{aligned}
\end{equation}
$0 < \al < 2$ かつ $\ga > 0$ ならば、$\it{x} \rightarrow z^\star$ が示され、ここで $z^\star$ は~\eqref{eq-splitting-opt}の解である。詳細は~\citep{Combettes2007}を参照されたい。
%
このアルゴリズムは、もう一つのよく知られた手法である交互方向乗数法~\citep{GabayMercier,GlowinskiMarroco}（レビューについては~\citep{BoydADMM}も参照）と密接に関連しており、これは双対問題にDRを適用することで得られる。二つの手法の等価性（最初に~\citep{Eckstein1992}によって示された）の詳細については~\citep{FPapPeyOud13}を参照されたい。

問題~\eqref{eq-bb-discrete}を形式~\eqref{eq-splitting-opt}に再定式化する方法は多数あり、これらのアプローチの詳細な説明については~\citep{FPapPeyOud13}を参照されたい。簡単な方法として、$x=(\a,\MomentD,\tilde\a,\tilde\MomentD)$ と設定し、
\eq{
	F(x) \eqdef \Theta(\tilde\a,\tilde\MomentD) + \iota_{\BBConstrD(\a_0,\a_1)}(\a,\MomentD)
	\qandq
	G(x) = \iota_{\Dd}(\a,\MomentD,\tilde\a,\tilde\MomentD),
}
\eq{
	\qwhereq \Dd \eqdef \enscond{(\a,\MomentD,\tilde\a,\tilde\MomentD)}{
		\tilde\a=\InterpBB_a(\a), \tilde\MomentD=\InterpBB_\Moment( \MomentD )
	}
}
とする。
%
これら二つの関数の近接作用素は効率的に計算できる。
%
実際、
\eq{
	\Prox_{\tau F}(x) = ( \Prox_{\tau\Theta}(\tilde\a,\tilde\MomentD), \Proj_{\BBConstrD(\a_0,\a_1)}(\a,\MomentD) )
}
が成り立つ。近接作用素 $\Prox_{\tau\Theta}$ は各格子点で三次多項式方程式を解くことで計算される。アフィン制約 $\BBConstrD(\a_0,\a_1)$ への直交射影はPoisson方程式の求解を含み、高速Fourier変換を用いて $O(N\log(N))$ 演算で達成できる。ここで $N=Tn_1 n_2$ は格子点の数である。
%
最後に、近接作用素 $\Prox_{\tau G}$ は線形射影であり、小さな線形系の逆行列計算を必要とする。
%
これらの計算の詳細については~\citet{FPapPeyOud13}を参照されたい。
%
図~\ref{fig-dynamic-maze}に、複雑な平面領域内の動的補間を計算するためにこの手法を用いた例を示す。
%
動的OTに対するこの種の近接法は、平均場ゲーム~\citep{benamou2015augmented}などの関連する問題を解くためにも使用されている。

\begin{figure}[h!]
\centering
\includegraphics[width=\linewidth]{dynamic/dynamic-maze}
\caption{\label{fig-dynamic-maze}
近接分離スキームを用いて計算された動的OTの解 $\al_t$。
}
\end{figure}



%%%%%%%%%%%%%%%%%%%%%%%%%%%%%%%%%%%%%%%%%
\section{動的非平衡OT}
\label{dynamic-unbalanced}

異なる質量 $\al_0(\Xx) \neq \al_1(\Xx)$ を持つ入力測度を対応付けるため（いわゆる「非平衡」設定、\citet{benamou2003numerical}によって導入された用語）、また局所的な質量変動に対処するため、
%
いくつかの正規化や緩和が提案されている。特に固定周辺分布制約を緩和する方法がある。\S\ref{sec-unbalanced}を参照されたい。
一般的な方法論は、連続の方程式~\eqref{eq-mass-conservation}にソース項 $s_t(x)$ を導入することからなる。したがって次を考える：
\eq{
	\bar\BBConstr(\al_0,\al_1) \eqdef
	\enscond{ (\al_t,\Moment_t,s_t) }{
		\pd{\al_t}{t} + \diverg(\Moment_t) = s_t,
		\al_{t=0}=\al_0,
		\al_{t=1}=\al_1
	}.
}
重要な問題は、このソース項に関連するコストをどのように測定し、元の動的定式化~\eqref{eq-bb-convex}に導入するかである。文献にはいくつかの提案があり、例えば $L^2$ コストを用いるもの~\citet{piccoli2014generalized}がある。
%
質量を「テレポート」すること（無限速度で移動し、それまで質量がなかった領域に突然質量が現れること）を避けるためには、関連するコストは無限大であるべきである。これは、$s_t$ も任意の測度であることを許容し（例えば1次斉次コストを用いて）、$s_t$ を運動量 $\Moment_t$ と同様にペナルティ化することで、単純な凸的方法で達成できることがわかる：
\eql{\label{eq-wfr-dynamic}
	\WFR^2(\al_0,\al_1) = \umin{ (\al_t,\Moment_t,s_t)_t  \in \bar\BBConstr(\al_0,\al_1)  }
		\Theta(\al,\Moment,s),
}
\eq{
	\qwhereq \Theta(\al,\Moment,s) \eqdef
		\int_0^1 \int_{\RR^d} \pa{
			\theta(\al_t(x),\Moment_t(x)) + \tau \theta(\al_t(x),s_t(x))
		}\d x \d t,
}
ここで $\th$ は~\eqref{eq-theta-dynamic}で導入された凸1次斉次関数であり、$\tau$ は質量の輸送と質量の生成/消滅間のトレードオフを制御する重みである。この定式化は複数の著者によって独立に提案された~\citep{LieroMielkeSavareShort,2017-chizat-focm,kondratyev2015}。この「動的」定式化には「静的」な対応物がある。注意~\ref{rem-wfr-static}を参照されたい。
%
凸最適化問題~\eqref{eq-wfr-dynamic}は、\S\ref{sec-prox-solvers}で詳述した方法と類似の手法を用いて解くことができる。
%
図~\ref{fig-unbalanced-dynamic}にいくつかの非平衡OT動的補間の比較を示す。
%
この動的定式化は形状レジストレーションの「メタモルフォーシス」モデル~\citep{Metamorphosis2005}に類似しており、より正確な関連は~\citep{maas2015generalized,maas2016generalized}に詳述されている。

$\tau \rightarrow 0$ かつ $\al_0(\Xx) = \al_1(\Xx)$ の場合、古典的なOT問題が回復され、$\WFR(\al_0,\al_1) \rightarrow \Wass(\al_0,\al_1)$ となる。一方、$\tau \rightarrow +\infty$ の場合、この距離は密度上のヘリンガー計量に近づく：
\begin{align*}
	\frac{1}{\tau}\WFR(\al_0,\al_1)^2\; \overset{\tau \rightarrow +\infty} \longrightarrow
	&\int_{\Xx} | \sqrt{\density{\al_0}(x)} - \sqrt{\density{\al_1}(x)} |^2 \d x \\
	=& \int_{\Xx} | 1 - \sqrt{ \frac{\d\al_1}{\d\al_0}(x) } |^2 \d \al_0(x).
\end{align*}

\begin{figure}[h!]
\centering
\includegraphics[width=\linewidth]{unbalanced/unbalanced-dynamic}
\caption{\label{fig-unbalanced-dynamic}
ヘリンガー（第1行）、Wasserstein（第2行）、部分最適輸送（第3行）、Wasserstein--Fisher--ラオ（第4行）の動的補間の比較。
}
\end{figure}




%%%%%%%%%%%%%%%%%%%%%%%%%%%%%%%%%%%%%%%%%
\section{より一般的な移動度汎関数}
\label{eq-generic-modibility}

動的定式化~\eqref{eq-bb-convex}を、\eqref{eq-theta-dynamic}で定義されたものの代わりに他の「移動度関数」$\th$ を考えることで一般化することが可能である。この移動度汎関数の可能な選択肢は~\citet{dolbeault2009new}によって提案されている：
\eql{\label{eq?generic-mobility}
	\foralls (a,b) \in \RR_+ \times \RR^\dim, \quad
	\th(a,b) = a^{s-p} \norm{b}^{p},
}
ここでパラメータは $\th$ が凸となるために $p \geq 1$ かつ $s \in [1,p]$ を満たす必要がある。
%
$1<s\leq p$ の場合、$\th$ は無限大で線形増大しないため、\eqref{eq-bb-convex}の解はLebesgue測度に関する密度を持つように制約される必要があり、この定義は注意して扱う必要があることに留意されたい。

$s=1$ の場合は古典的なOT問題に対応し、\eqref{eq-bb-convex}の最適値は $\Wass_p(\al,\be)$ を定義する。この場合、$\th$ は1次斉次であり、\eqref{eq-bb-convex}の解は任意の測度であり得る。$(s=1,p=2)$ の場合は $\Wass_2$ を定義するために~\eqref{eq-bb-convex}で考慮された初期設定である。

極限的な場合 $s=p$ も興味深い。これは双対Sobolevノルム $W^{-1,p}$ に対応し、\eqref{eq-bb-convex}の値は
\eq{
	\norm{\al-\be}_{W^{-1,p}(\RR^\dim)}^p =
	\umin{\f} \enscond{ \int_{\RR^\dim} \f  \d(\al-\be)}{  \int_{\RR^\dim} \norm{ \nabla f(x) }^q \d x \leq 1 }
}
に等しくなる。ここで $1/q+1/p=1$ である。
%
極限 $(p=s,q) \rightarrow (1,\infty)$ では、$\Wass_1$ ノルムが回復される。
%
$s=p=2$ の場合は~\eqref{eq-dual-sobolev-div}で定義されたSobolev $H^{-1}(\RR^\dim)$ Hilbertノルムに対応する。


\todoK{Speak of the connexion with geodesic in shape space~\citet{beg2005computing}. }

\todoK{``Here many computation techniques related to optimal control and Hamilton-Jacobi equations have not mentioned. The connections between optimal transport and Mean field games may be the other interesting direction. Many techniques developed in numerical optimal transport can be applied to general control problems in the set of probability space. For example, a new Hopf-Lax formula is introduced for dynamical optimal transport problem. See related discussions in~\citet{gangbo2008hamilton}.''}



%%%%%%%%%%%%%%%%%%%%%%%%%%%%%%%%%%%%%%%%%
\section{経路空間上の動的定式化}
\label{sec-entropic-dynamic}

古典的OTおよびエントロピー正則化された（\S\ref{c-entropic}を参照）OTの定式化には、空間 $\Xx$ 上のすべての可能な経路 $\ga : [0,1] \rightarrow \Xx$（すなわち曲線）の空間 $\bar\Xx$ 上の抽象的な最適化問題を研究することに基づく、自然な動的定式化が存在する。
%
簡単のため $\Xx=\RR^\dim$ と仮定するが、これは測地空間やグラフなどのより一般的な空間に拡張される。
%
非形式的には、時刻 $t=0,1$ における二つの入力測度 $\al_0,\al_1$ 間の「粒子」の動的挙動は、確率分布 $\bar\pi \in \Mm_+^1(\bar\Xx)$ によって記述される。このような分布は、始点と終点の分布が $(\al_0,\al_1)$ と一致しなければならないという条件を満たすべきであり、これは押し出しを用いて形式的に次のように書かれる：
\eq{
	\bar\Couplings(\al_0,\al_1) \eqdef \enscond{ \bar\pi \in \Mm_+^1(\bar\Xx) }{ \bar P_{0\sharp} \bar\pi = \al_0, \bar P_{1\sharp} \bar\pi = \al_1  },
}
ここで任意の経路 $\ga \in \bar\Xx$ に対して $P_0(\ga) = \ga(0), P_1(\ga) = \ga(1)$ である。

%%%%
\paragraph{経路空間上のOT}

経路空間上で定式化された古典的OT~\eqref{eq-mk-generic}の動的版は次のようになる：
\eql{\label{eq-ot-pathsspace}
	\Wass_2(\al_0,\al_1)^2 =
	\umin{ \bar\pi \in \bar\Couplings(\al_0,\al_1) } \int_{\bar\Xx} \Ll(\ga)^2 \d\bar\pi(\ga),
}
ここで $\Ll(\ga) = \int_0^1 |\ga'(s)|^2 \d s$ は経路 $s \in [0,1] \mapsto \ga(s) \in \Xx$ の運動エネルギーである。それぞれ~\eqref{eq-ot-pathsspace}と~\eqref{eq-mk-generic}を解く最適カップリング $\pi^\star$ と $\bar\pi^\star$ の関連は、$\bar\pi^\star$ が $\pi^\star$ によって規定された割合で点の対を結ぶ測地線にのみ質量を与えることである。離散測度の特別な場合、これは
\eq{
	\pi^\star = \sum_{i,j} \P_{i,j} \de_{(x_i,y_j)}
	\qandq
	\bar\pi^\star = \sum_{i,j} \P_{i,j} \de_{\ga_{x_i,y_j}},
}
を意味する。ここで $\ga_{x_i,y_j}$ は $x_i$ と $y_j$ 間の測地線である。さらに、$\ga$ が $\bar\pi^\star$ に従って抽出されるとき、時刻 $t$ における曲線点 $\ga(t)$ の分布によって定義される測度、すなわち
\eql{\label{eq-interpolation-spacepaths}
	 t \in [0,1] \mapsto \al_t \eqdef P_{t\sharp} \bar\pi^\star
	 \qwhereq
	 P_t(\ga) = \ga(t) \in \Xx,
}
は動的定式化~\eqref{eq-bb-convex}の解、すなわち変位補間である。
%
離散的な場合には、\eqref{eq-mccann-discrete}が回復される。

%%%%
\paragraph{経路空間上のエントロピーOT}

次に、第~\ref{c-entropic}章で定義されたエントロピーOTの、経路空間を用いた再解釈に移る。
%
\eqref{eq-entropic-generic-proj}と同様に、これはKullback--Leiblerダイバージェンスの射影を用いて定義されるが、今回は経路空間上の参照測度 $\bar\Kk$（可逆Brown運動（Wiener過程）の分布であり、初期時刻と最終時刻で一様分布を持つ）に対するものである：
\eql{\label{eq-shrodinger-paths}
	\umin{ \bar\pi \in \bar\Couplings(\al_0,\al_1) } \KL(\bar\pi|\bar\Kk).
}
この問題の概観と~\citet{Schroedinger31}の仕事に関する歴史的説明については、\citet{LeonardSchroedinger}のレビュー論文を参照されたい。
%
\eqref{eq-shrodinger-paths}の（一意的な）解 $\bar\pi_\varepsilon^\star$ は $\varepsilon \rightarrow 0$ のとき~\eqref{eq-ot-pathsspace}の解に収束することが示される。さらに、この解はBrown橋 $\bar\ga_{x,y}^\varepsilon \in \bar\Xx$（ファジーな測地線に類似し、$\varepsilon \rightarrow 0$ で $\de_{\ga_{x,y}}$ に収束する）を用いて、静的エントロピーOT問題~\eqref{eq-entropic-generic}の解と結びつけられる。離散的な設定では、これは
\eql{\label{eq-brownian-bridge-interp}
	\pi_\varepsilon^\star = \sum_{i,j} \P_{\varepsilon,i,j}^\star \de_{(x_i,y_j)}
	\qandq
	\bar\pi_\varepsilon^\star = \sum_{i,j} \P_{\varepsilon,i,j}^\star \bar\ga_{x_i,y_j}^\varepsilon,
}
を意味する。ここで $\P_{\varepsilon,i,j}^\star$ はSinkhornアルゴリズムを用いて計算できる。
%
\eqref{eq-interpolation-spacepaths}と同様に、エントロピー補間を次のように定義できる：
\eq{
	\al_{\varepsilon,t} \eqdef \P_{t\sharp} \bar\pi_\varepsilon^\star.
}
Brown橋に沿った時刻 $t$ における位置の分布 $\P_{t\sharp} \bar\ga_{x,y}^\varepsilon$ は、$\ga_{x,y}(t)$ を中心とする分散 $t(1-t)\varepsilon^2$ のGauss分布 $\Gg_{t(1-t)\varepsilon^2}(\cdot-\ga_{x,y}(t))$ であるから、$\al_{\varepsilon,t}$ は移動するDirac測度の集合のGaussぼかしであることが導かれる：
\eq{
	\al_{\varepsilon,t} =
	\sum_{i,j} \P_{\varepsilon,i,j}^\star \Gg_{t(1-t)\varepsilon^2}( \cdot - \ga_{x_i,y_j}(t) ).
}
得られたBrown橋の混合は図~\ref{fig-schrodinger-dynamic}に表示されている。


% G B D H
\newcommand{\MyFigSchrDyn}[1]{\includegraphics[width=.245\linewidth,trim=50 45 38 35,clip]{schrodinger-dynamic/schrodinger-dynamic-eps#1}}
%%% FIG %%%
\begin{figure}[h!]\centering
\centering
\begin{tabular}{@{}c@{}c@{}c@{}c@{}}
\MyFigSchrDyn{5}&
\MyFigSchrDyn{50}&
\MyFigSchrDyn{200}&
\MyFigSchrDyn{1000}\\
$\varepsilon=0$ & $\varepsilon=.05$ & $\varepsilon=0.2$ & $\varepsilon=1$
\end{tabular}
%%%
\caption{%%%
Schrödingerのエントロピー補間~\eqref{eq-brownian-bridge-interp}に付随する経路空間上のBrown橋経路のサンプル。青は $t=0$ に、赤は $t=1$ に対応する。
} \label{fig-schrodinger-dynamic}
\end{figure}
%%% FIG %%%

このエントロピー補間 $(\al_t)_t$ を記述するもう一つの方法は、Benamou--Brenierの動的定式化~\eqref{eq-bb-non-cvx}の正則化を用いるものであり、すなわち \todoK{check how $\varepsilon$ enters the picture}
\eql{\label{eq-bb-non-cvx-entropy}
	\umin{ (\al_t,\speed_t)_t \text{ sat. \eqref{eq-cons-bb-ncvx}} }
		\int_0^1 \int_{\RR^d}  \pa{
			 \norm{\speed_t(x)}^2  + \frac{\varepsilon}{4} \norm{ \nabla \log(\al_t)(x) }^2
			 }
			 \d\al_t(x) \d t;
}
\citep{gentil2015analogy,chen2016relation}を参照されたい。


